% Options for packages loaded elsewhere
\PassOptionsToPackage{unicode}{hyperref}
\PassOptionsToPackage{hyphens}{url}
\documentclass[
  11pt,
]{article}
\usepackage{xcolor}
\usepackage[left=3cm,right=3cm,top=2.5cm,bottom=2.5cm]{geometry}
\usepackage{amsmath,amssymb}
\setcounter{secnumdepth}{5}
\usepackage{iftex}
\ifPDFTeX
  \usepackage[T1]{fontenc}
  \usepackage[utf8]{inputenc}
  \usepackage{textcomp} % provide euro and other symbols
\else % if luatex or xetex
  \usepackage{unicode-math} % this also loads fontspec
  \defaultfontfeatures{Scale=MatchLowercase}
  \defaultfontfeatures[\rmfamily]{Ligatures=TeX,Scale=1}
\fi
\usepackage{lmodern}
\ifPDFTeX\else
  % xetex/luatex font selection
\fi
% Use upquote if available, for straight quotes in verbatim environments
\IfFileExists{upquote.sty}{\usepackage{upquote}}{}
\IfFileExists{microtype.sty}{% use microtype if available
  \usepackage[]{microtype}
  \UseMicrotypeSet[protrusion]{basicmath} % disable protrusion for tt fonts
}{}
\usepackage{setspace}
\makeatletter
\@ifundefined{KOMAClassName}{% if non-KOMA class
  \IfFileExists{parskip.sty}{%
    \usepackage{parskip}
  }{% else
    \setlength{\parindent}{0pt}
    \setlength{\parskip}{6pt plus 2pt minus 1pt}}
}{% if KOMA class
  \KOMAoptions{parskip=half}}
\makeatother
\usepackage{longtable,booktabs,array}
\usepackage{caption}
\captionsetup[table]{skip=6pt}
\newcounter{none} % for unnumbered tables
\usepackage{calc} % for calculating minipage widths
% Correct order of tables after \paragraph or \subparagraph
\usepackage{etoolbox}
\makeatletter
\patchcmd\longtable{\par}{\if@noskipsec\mbox{}\fi\par}{}{}
\makeatother
% Allow footnotes in longtable head/foot
\IfFileExists{footnotehyper.sty}{\usepackage{footnotehyper}}{\usepackage{footnote}}
\makesavenoteenv{longtable}
\usepackage{graphicx}
\makeatletter
\newsavebox\pandoc@box
\newcommand*\pandocbounded[1]{% scales image to fit in text height/width
  \sbox\pandoc@box{#1}%
  \Gscale@div\@tempa{\textheight}{\dimexpr\ht\pandoc@box+\dp\pandoc@box\relax}%
  \Gscale@div\@tempb{\linewidth}{\wd\pandoc@box}%
  \ifdim\@tempb\p@<\@tempa\p@\let\@tempa\@tempb\fi% select the smaller of both
  \ifdim\@tempa\p@<\p@\scalebox{\@tempa}{\usebox\pandoc@box}%
  \else\usebox{\pandoc@box}%
  \fi%
}
% Set default figure placement to htbp
\def\fps@figure{htbp}
\makeatother
% definitions for citeproc citations
\NewDocumentCommand\citeproctext{}{}
\NewDocumentCommand\citeproc{mm}{%
  \begingroup\def\citeproctext{#2}\cite{#1}\endgroup}
\makeatletter
 % allow citations to break across lines
 \let\@cite@ofmt\@firstofone
 % avoid brackets around text for \cite:
 \def\@biblabel#1{}
 \def\@cite#1#2{{#1\if@tempswa , #2\fi}}
\makeatother
\newlength{\cslhangindent}
\setlength{\cslhangindent}{1.5em}
\newlength{\csllabelwidth}
\setlength{\csllabelwidth}{3em}
\newenvironment{CSLReferences}[2] % #1 hanging-indent, #2 entry-spacing
 {\begin{list}{}{%
  \setlength{\itemindent}{0pt}
  \setlength{\leftmargin}{0pt}
  \setlength{\parsep}{0pt}
  % turn on hanging indent if param 1 is 1
  \ifodd #1
   \setlength{\leftmargin}{\cslhangindent}
   \setlength{\itemindent}{-1\cslhangindent}
  \fi
  % set entry spacing
  \setlength{\itemsep}{#2\baselineskip}}}
 {\end{list}}
\usepackage{calc}
\newcommand{\CSLBlock}[1]{\hfill\break\parbox[t]{\linewidth}{\strut\ignorespaces#1\strut}}
\newcommand{\CSLLeftMargin}[1]{\parbox[t]{\csllabelwidth}{\strut#1\strut}}
\newcommand{\CSLRightInline}[1]{\parbox[t]{\linewidth - \csllabelwidth}{\strut#1\strut}}
\newcommand{\CSLIndent}[1]{\hspace{\cslhangindent}#1}
\setlength{\emergencystretch}{3em} % prevent overfull lines
\providecommand{\tightlist}{%
  \setlength{\itemsep}{0pt}\setlength{\parskip}{0pt}}
% Paragraph scaffolding for manuscript revision.
%
% Two environments, presented in this order per paragraph:
%
%   rgt   -- author's rewrite (initially a placeholder)
%            small, italic, dark blue
%   orig  -- original Claude-authored draft, and until the rewrite
%            replaces it, the body text of the manuscript
%            normal size, gray
%
% The `bullets' environment (a boiled-down topic sentence plus bullet
% points, one per paragraph) is no longer used inside manuscripts:
% those blocks now live in a companion bullets.Rmd supplement in the
% same report directory, which typesets them as plain \textbf{} +
% itemize and needs no part of this file. The definition is retained
% only so that a stray block left in a document still compiles.
%
% This file is identical across every manuscript that includes it.
% It previously existed in two variants that swapped the sizes of
% rgt and orig, so the same markup rendered differently depending on
% which repository the manuscript lived in. The variant retained here
% keeps orig at full size because orig, not rgt, currently carries
% the prose in every manuscript in the pool.

\usepackage{xcolor}

\newcounter{pgraph}

\newenvironment{bullets}{%
  \stepcounter{pgraph}%
  \begin{quote}\small\sffamily\color[gray]{0.30}%
  {\normalfont\scriptsize\color[gray]{0.58}[\thepgraph]\enspace}}{%
  \end{quote}}

\newenvironment{rgt}{%
  \begin{quote}\small\itshape\color[RGB]{20,60,160}}{%
  \end{quote}}

\newenvironment{orig}{%
  \begin{quote}\normalsize\color[gray]{0.55}}{%
  \end{quote}}

% backward compat alias
\newenvironment{claudecode}{\begin{orig}}{\end{orig}}
%% tools/stamp.tex
%% Provenance footer for PDFs rendered from this project.
%% Vendored by zzcollab; this file is not project-specific. The
%% footer prints the source document and the time of rendering.
%% The git version is defined here too, and so is available to a
%% project that wants it on the page, but it is left off the
%% default footer: it already names the staged copy in share/ and
%% appears in its manifest row. All values are supplied at render
%% time by tools/stamp-render.R, which writes a small generated
%% file that redefines the macros below. The \providecommand
%% defaults keep a bare LaTeX compile working when the document is
%% built without the wrapper.
\usepackage{fancyhdr}
\usepackage{xcolor}
\providecommand{\stampsource}{\detokenize{(source not recorded)}}
\providecommand{\stamptime}{(time not recorded)}
\providecommand{\stampversion}{\detokenize{(version not recorded)}}
\setlength{\footskip}{40pt}
\newcommand{\stampfootcontent}{%
  \footnotesize\color{gray}%
  \parbox[b]{\textwidth}{%
    \stamptime\hfill\thepage\\[2pt]%
    \ttfamily\raggedright\stampsource}}
\pagestyle{fancy}
\fancyhf{}
\fancyfoot[C]{\stampfootcontent}
\renewcommand{\headrulewidth}{0pt}
\renewcommand{\footrulewidth}{0.4pt}
%% The title page uses the `plain` page style; redefine it so the
%% stamp appears there too.
\fancypagestyle{plain}{%
  \fancyhf{}%
  \fancyfoot[C]{\stampfootcontent}%
  \renewcommand{\headrulewidth}{0pt}%
  \renewcommand{\footrulewidth}{0.4pt}}
\renewcommand{\stampsource}{\textasciitilde{}/\allowbreak{}prj/\allowbreak{}res/\allowbreak{}04-runin-power-analysis/\allowbreak{}runinpower/\allowbreak{}analysis/\allowbreak{}report/\allowbreak{}03-allocation/\allowbreak{}report.Rmd}
\renewcommand{\stamptime}{2026-08-15 16:32 PDT}
\renewcommand{\stampversion}{\detokenize{455c889-wip}}
\usepackage{amsmath}
\usepackage{amssymb}
\usepackage{booktabs}
\usepackage{longtable}
\usepackage{float}
\usepackage[mathlines]{lineno}
\linenumbers
\DeclareMathOperator{\Var}{Var}
\DeclareMathOperator{\Cov}{Cov}
\DeclareMathOperator{\Corr}{Corr}
\usepackage{bookmark}
\IfFileExists{xurl.sty}{\usepackage{xurl}}{} % add URL line breaks if available
\urlstyle{same}
% fallback for those not using the hyperref driver hyperxmp:
\makeatletter
\@ifundefined{xmpquote}{\newcommand{\xmpquote}[1]{#1}}{}
\makeatother
\hypersetup{
  pdftitle={Optimal Allocation of Observations across Run-In{,} Treatment{,} and Common Close Phases in Longitudinal Clinical Trials},
  pdfauthor={Ronald G. Thomas, Ph.D.},
  hidelinks,
  pdfcreator={LaTeX via pandoc}}

\title{Optimal Allocation of Observations across Run-In,
Treatment, and Common Close Phases in Longitudinal
Clinical Trials}
\author{Ronald G. Thomas, Ph.D.\thanks{Wertheim School of Public Health, UCSD; ORCID: 0000-0003-1686-4965}}
\date{August 15, 2026}

\begin{document}
\maketitle
\begin{abstract}
\textbf{Background.} Trial designers planning a longitudinal
study with a fixed total observation budget face an
allocation question: how should observations be distributed
across run-in (\(J_0\)), treatment (\(J_1\)), and common close
(\(J_2\)) phases to minimize the variance of the
treatment-effect estimator? Existing guidance is heuristic;
the optimum has not been characterized for the three-phase
run-in setting.
\textbf{Methods.} We formulate the allocation question as a
discrete optimization problem over the Woodbury-based
variance formula derived in Paper 1 of this compendium and
solve it by exhaustive enumeration of all feasible phase
splits. The behavior of the optimum is described as a
function of the signal-to-noise ratio \(\sigma_b / \sigma\).
Extensions cover replicated endpoint observations (closely
spaced measurements at the end of treatment that average
out measurement error on the final assessment, analogous to
run-in averaging applied at the post-treatment end) and
pattern-mixture handling of monotone dropout following the
Frost-Kenward-Fox (2008) framework.
\textbf{Results.} Under a fixed total budget the optimal design
places zero observations in the run-in phase in every
configuration examined; the budget is split between the
treatment and common close phases. This reconciles with the
run-in efficiency gains of Paper 1: run-in observations help
when they are additional visits appended to a fixed
treatment phase, but when visits are fungible the run-in
phase is never chosen, because it introduces the
pre-randomization slope parameter without contributing to
the treatment contrast. The common close phase enters the
optimal design once \(\sigma_b / \sigma\) exceeds a
budget-dependent crossover, which falls from 0.81 at \(J=3\)
to 0.30 at \(J=8\); at high ratios roughly half the budget
goes to the common close. Replicating the final endpoint
reduces the treatment-effect variance by up to 31\% at
\(\sigma_b / \sigma = 0.5\) (five replicates) but by at most
5\% at \(\sigma_b / \sigma = 2\).
\textbf{Conclusions.} The numerical allocation results give
trial designers a simple decision procedure driven by the
signal-to-noise ratio and the visit budget. The procedure
extends to replicated-endpoint and dropout-adjusted designs
within the same variance framework.
\end{abstract}

{
\setcounter{tocdepth}{2}
\tableofcontents
}
\setstretch{1.1}
\section{Introduction}\label{introduction}

\begin{orig}
The relative efficiency surface in Paper 1 (Figure 1)
demonstrates diminishing returns from additional
run-in observations appended to a fixed treatment
phase: most of the attainable variance reduction is
recovered by the first two run-in observations, and
the marginal value falls quickly thereafter (Section
3 examines the marginal pattern, including a penalty
specific to a single run-in observation).
Similarly, common close observations improve efficiency
but at a rate that depends on the number of treatment
observations \(J_1\).

\end{orig}

\begin{orig}
These observations raise a design question that Paper 1
does not address: given operational constraints on the
total number of visits or total trial duration, what is
the optimal distribution of observations across the
three phases?

\end{orig}

\begin{orig}
\textsuperscript{\citeproc{ref-winkens2005}{1}} calculated optimal time-point positions
for linearly divergent treatment effects under compound
symmetry, but did not consider the run-in setting.\textsuperscript{\citeproc{ref-yi2002}{2}} derived sample sizes for trends across repeated
measurements under missing data.\textsuperscript{\citeproc{ref-frison1992}{3}} showed that having more than one
pre-treatment measurement is beneficial for designs
analyzed via ANCOVA.\textsuperscript{\citeproc{ref-overall1994}{4}} provided sample size
formulas for repeated measurement designs.\textsuperscript{\citeproc{ref-galbraith2002}{5}} developed general guidelines for the
design of clinical trials with longitudinal outcomes.

The question posed here is a design-optimization problem
and belongs to a well-developed tradition. The general
theory of optimum experimental design, in which a design
is chosen to optimize a functional of the Fisher
information matrix, is set out by\textsuperscript{\citeproc{ref-atkinson1992}{6}} and\textsuperscript{\citeproc{ref-fedorov1997}{7}}.\textsuperscript{\citeproc{ref-mentre1997}{8}} extended the theory to
random-effects regression models, the model class used
here, giving an algorithm that maximizes the determinant
of the population information matrix. Closest in setting
is\textsuperscript{\citeproc{ref-tekle2008}{9}}, who constructed D-optimal cohort designs
for linear mixed-effects models with a cost function for
repeated measurements, and found that concentrating
measurements at optimally chosen time points in a single
cohort is efficient. Three features distinguish the
present problem from that literature. First, the design
space is not a set of time points to be placed freely
but a partition of a fixed count of equally spaced
visits into three phases with a prescribed ordering.
Second, the run-in phase introduces its own nuisance
parameter, the pre-randomization slope \(\delta\), so
enlarging that phase both adds information and consumes
a degree of freedom; the optimal-design literature
usually holds the parameter vector fixed as the design
varies. Third, the criterion here is the variance of a
single scalar contrast, a \(c\)-optimality criterion,
rather than the D-optimality that most of this
literature adopts. The results below are accordingly
numerical rather than an application of the general
equivalence theory.

The slopepower command\textsuperscript{\citeproc{ref-nash2021}{10}} and the longpower
R package\textsuperscript{\citeproc{ref-iddi2022}{11}} implement power calculations
for slope-comparison trials but do not optimize the
allocation of observations across phases.

\end{orig}

\section{Problem formulation}\label{problem-formulation}

\subsection{Fixed total visits}\label{fixed-total-visits}

\begin{orig}
Given \(J\) total observations (excluding baseline),
choose \((J_0^*, J_1^*, J_2^*)\) to minimize
\(\mathop{\mathrm{Var}}_1(\hat\gamma)\) subject to \(J_0 + J_1 + J_2 = J\),
\(J_0 \ge 0\), \(J_1 \ge 1\), \(J_2 \ge 0\).

The objective function is
\(V(J_0, J_1, J_2) = [(X'\Sigma^{-1}X)^{-1}]_{\gamma\gamma}\)
computed via the Woodbury identity from Paper 1.

\end{orig}

Throughout, the phase counts \((J_0, J_1, J_2)\) correspond
to the arguments \((r, k, f)\) of the \texttt{runinpower} package
routines used in the code chunks and tables.

\subsection{Fixed total duration}\label{fixed-total-duration}

\begin{orig}
Alternatively, fix the trial duration
\(T = (J_0 + J_1 + J_2) \cdot t\) and minimize
\(\mathop{\mathrm{Var}}_1(\hat\gamma)\) subject to
\(J_0 + J_1 + J_2 = T/t\), or more generally allow
non-equally-spaced observations within each phase
and minimize over the observation times as well.

\end{orig}

\section{Diminishing returns from run-in observations}\label{diminishing-returns-from-run-in-observations}

\subsection{\texorpdfstring{Marginal value of the \(J_0\)-th run-in observation}{Marginal value of the J\_0-th run-in observation}}\label{marginal-value-of-the-j_0-th-run-in-observation}

Define the marginal variance reduction from adding
the \(J_0\)-th run-in observation:
\begin{equation}
  \Delta V_{J_0} = V(J_0 - 1, J_1, J_2) - V(J_0, J_1, J_2)
\end{equation}

\begin{table}

\caption{\label{tab:marginal-value}Marginal variance reduction from the $J_0$-th run-in observation ($J_1=2$, $J_2=0$, $\sigma^2 = 1$, $\sigma_b^2 = 1$).}
\centering
\begin{tabular}[t]{rr}
\toprule
$J_0$ & $\Delta V_{J_0}$\\
\midrule
1 & 0.1765\\
2 & 0.6667\\
3 & 0.5124\\
4 & 0.3176\\
5 & 0.2008\\
\addlinespace
6 & 0.1342\\
7 & 0.0945\\
8 & 0.0696\\
\bottomrule
\end{tabular}
\end{table}

\subsection{\texorpdfstring{The \(J_0=1\) anomaly}{The J\_0=1 anomaly}}\label{the-j_01-anomaly}

\begin{orig}
A non-monotonicity is visible in the table:
\(\Delta V_1 < \Delta V_2\) at the displayed parameter
setting (\(\sigma_b / \sigma = 1\)). The first run-in
observation is less valuable than the second. This
occurs because adding a single run-in observation
forces the model from 2 parameters (\(\beta, \gamma\))
to 3 parameters (\(\delta, \beta, \gamma\)), consuming a
degree of freedom. The second run-in observation does
not increase the parameter count and provides pure
information gain. The reversal is not universal: a
numerical scan (Table \ref{tab:anomaly-boundary})
shows that, for \(J_1 = 2\), \(J_2 = 0\), and
\(\sigma = t = 1\), the inequality
\(\Delta V_1 < \Delta V_2\) holds only for
\(\sigma_b / \sigma\) in an intermediate band,
approximately \((0.53, 1.51)\); at lower and higher
ratios the first run-in observation is the more
valuable one.

\end{orig}

\begin{table}

\caption{\label{tab:anomaly-boundary}Boundary of the $J_0 = 1$ anomaly ($J_1 = 2$, $J_2 = 0$, $\sigma = t = 1$). The reversal $\Delta V_1 < \Delta V_2$ holds only between the two computed boundary ratios 0.53 and 1.51.}
\centering
\begin{tabular}[t]{rrrl}
\toprule
$\sigma_b/\sigma$ & $\Delta V_1$ & $\Delta V_2$ & $\Delta V_1 < \Delta V_2$\\
\midrule
0.32 & 0.1633 & 0.0080 & no\\
0.53 & 0.0454 & 0.0418 & no\\
0.75 & 0.0033 & 0.2653 & yes\\
1.00 & 0.1765 & 0.6667 & yes\\
1.25 & 0.6649 & 1.1082 & yes\\
\addlinespace
1.51 & 1.5407 & 1.5352 & no\\
2.00 & 4.2000 & 2.1718 & no\\
\bottomrule
\end{tabular}
\end{table}

\begin{orig}
This effect is loosely analogous to\textsuperscript{\citeproc{ref-frison1992}{3}}'s finding
that a single pre-treatment measurement used as an ANCOVA
covariate can reduce power when the pre-post correlation
is low. The practical implication is regime-specific:
\textbf{within the intermediate signal-to-noise band of Table
\ref{tab:anomaly-boundary}, a trial designer who adds
run-in observations should plan at least two, not one.}
Outside that band the first run-in observation carries
the larger marginal value and no special caution applies.

\end{orig}

\subsection{Correlation structure governs diminishing returns}\label{correlation-structure-governs-diminishing-returns}

The correlation between the \(J_0\)-th run-in change
score and the treatment-period change scores
determines the information gain. Under compound
symmetry:
\begin{equation}
  \mathop{\mathrm{Corr}}(C_{-J_0}, C_l) = \frac{\sigma_b^2 t^2}
    {2\sigma^2 + \sigma_b^2 t^2}
\end{equation}

\begin{orig}
which is constant in \(J_0\). The diminishing returns
(for \(J_0 \ge 2\)) arise from the information matrix
structure, not from declining correlation. Under AR(1)\textsuperscript{\citeproc{ref-winkens2007}{12}}, the correlation does decline with
temporal distance, producing faster diminishing
returns from distant run-in observations.

\end{orig}

\section{\texorpdfstring{Optimal allocation for fixed \(J\)}{Optimal allocation for fixed J}}\label{optimal-allocation-for-fixed-j}

\begin{table}

\caption{\label{tab:optimal-grid}Optimal allocation $(J_0^*, J_1^*, J_2^*)$ for fixed total observations $J$, by signal-to-noise ratio.}
\centering
\begin{tabular}[t]{rllll}
\toprule
$J$ & Low ($\sigma_b/\sigma=0.32$) & Med ($\sigma_b/\sigma=1.0$) & High ($\sigma_b/\sigma=2.24$) & AD ($\sigma_b/\sigma=1.70$)\\
\midrule
3 & (0,3,0) & (0,2,1) & (0,2,1) & (0,2,1)\\
4 & (0,4,0) & (0,2,2) & (0,2,2) & (0,2,2)\\
5 & (0,5,0) & (0,3,2) & (0,3,2) & (0,3,2)\\
6 & (0,6,0) & (0,3,3) & (0,3,3) & (0,3,3)\\
7 & (0,5,2) & (0,4,3) & (0,4,3) & (0,4,3)\\
\addlinespace
8 & (0,5,3) & (0,4,4) & (0,4,4) & (0,4,4)\\
\bottomrule
\end{tabular}
\end{table}

\begin{orig}
The dominant feature of Table \ref{tab:optimal-grid} is
that the optimal run-in count is zero in every cell:
across all budgets \(J = 3, \ldots, 8\) and all four
signal-to-noise settings, including the Alzheimer's
neuroimaging regime, the optimizer splits the budget
between the treatment and common close phases and never
assigns a visit to the run-in phase. The mechanism is a
degree-of-freedom cost. Any positive run-in count forces
the model to estimate the pre-randomization slope
\(\delta\) in addition to \((\beta, \gamma)\), and run-in
observations inform the treatment contrast only
indirectly, through the random slope; a visit placed in
the treatment or common close phase buys the same slope
information without the extra parameter. This finding
does not contradict the efficiency gains of Paper 1.
There, run-in observations are additional visits appended
to a design whose treatment phase is held fixed, and the
comparison is against the shorter design without them.
Here visits are fungible under a fixed total budget, and
a run-in visit always has a better use. The two framings
answer different design questions, and the distinction is
the central practical message of this paper.

\end{orig}

\section{When is the common close phase beneficial?}\label{when-is-the-common-close-phase-beneficial}

\begin{orig}
The common close phase contributes information about
the random slope \(b_i\) without contributing to the
treatment comparison directly. When a common close
observation is added to an otherwise fixed design, it
reduced \(\mathop{\mathrm{Var}}(\hat\gamma)\) in every configuration we
examined (\(J_1 = 1, \ldots, 4\) with two run-in
observations, and \(\sigma_b/\sigma\) down to \(10^{-3}\)):
a free observation is always at least weakly helpful.
The design-relevant question is therefore not whether
the common close helps in isolation, but whether a
visit taken from a fixed budget is better spent on the
common close than on the treatment phase. That
trade-off is quantified next.

\end{orig}

\begin{table}

\caption{\label{tab:common-close-crossover}Signal-to-noise ratio above which reallocating the last budgeted visit from the treatment phase to the common close, $(0, J-1, 1)$ versus $(0, J, 0)$, reduces $\mathop{\mathrm{Var}}(\hat\gamma)$.}
\centering
\begin{tabular}[t]{rr}
\toprule
$J$ & Crossover $\sigma_b/\sigma$\\
\midrule
3 & 0.806\\
4 & 0.592\\
5 & 0.471\\
6 & 0.392\\
7 & 0.336\\
\addlinespace
8 & 0.295\\
\bottomrule
\end{tabular}
\end{table}

\subsection{\texorpdfstring{Crossover as a function of the budget \(J\)}{Crossover as a function of the budget J}}\label{crossover-as-a-function-of-the-budget-j}

\begin{orig}
Table \ref{tab:common-close-crossover} shows the
crossover ratio at which the reallocation of a single
budgeted visit from the treatment phase to the common
close becomes favorable. The crossover falls
monotonically with the budget, from
\(\sigma_b/\sigma \approx 0.81\) at \(J = 3\) to
\(\approx 0.30\) at \(J = 8\): larger trials should commit
visits to the common close at progressively lower
signal-to-noise ratios. Above the crossover the full
optimizer (Table \ref{tab:optimal-grid}) assigns an
increasing share of the budget to the close phase; at
\(J = 8\) the optimum is an even split \((0, 4, 4)\) at
every ratio examined from \(0.6\) upward, and even at
\(\sigma_b/\sigma = 0.32\) the optimum is \((0, 5, 3)\).
These values are numerical results from exhaustive
enumeration; an analytical characterization of the
crossover remains open.

\end{orig}

\section{Non-equally-spaced designs}\label{non-equally-spaced-designs}

\subsection{General time points}\label{general-time-points}

The display in this section comes from a Mathematica
notebook derivation for a \emph{different} model than the one
used elsewhere in this paper: a raw-outcome
random-slope-only model with centered observation times
(\(\sum t_i = 0\)) and independent residuals
\(R = \sigma^2 I\), rather than the change-score model
with shared baseline error
\(R = \sigma^2 (I + \mathbf{1}\mathbf{1}')\) that
underlies all preceding sections. Under that
raw-outcome model, with the normalization
\(\sum t_i^2 = 1\), the notebook gives for two time
points \(t_1, t_2\):
\begin{equation}
  \mathop{\mathrm{Var}}(\hat\gamma) = \frac{3\sigma^2}{1 - t_1 t_2}
    + 2\sigma_b^2
\end{equation}
This expression has not been re-derived or numerically
verified within the change-score framework of this
paper, and reconciling the two formulations remains
open.

\begin{orig}
Within the raw-outcome model above, the expression is
minimized when \(t_1 t_2\) is minimized (most negative),
i.e., when the two time points are maximally separated.
This is consistent with the known result that
observations at the extremes of the follow-up period
are optimal for slope estimation. Whether and how this
carries over to the change-score model of the preceding
sections has not been established here.

\end{orig}

\subsection{Optimization over observation times within phases}\label{optimization-over-observation-times-within-phases}

\section{Replicated endpoint observations}\label{replicated-endpoint-observations}

\subsection{Motivation}\label{motivation}

\begin{orig}
The run-in mechanism averages multiple pre-randomization
observations to reduce the contribution of measurement
error \(\sigma^2\) to the estimated baseline slope\textsuperscript{\citeproc{ref-frison1992}{3},\citeproc{ref-frost2008}{13}}.\textsuperscript{\citeproc{ref-goulet2019}{14}} showed that
replicated measures increase power but with diminishing
returns past 20-50 replications. The same logic applies
at the end of the treatment period:
if the final scheduled visit is replaced by (or augmented
with) two or three closely-spaced measurements, all
under the same treatment condition, the measurement
error on the endpoint assessment is reduced by
averaging.

\end{orig}

\begin{orig}
Concretely, if the last visit at time \(J_1 t\) is replicated
\(n_{\text{rep}}\) times over a short interval (days to
weeks, negligible relative to \(t\)), the effective
residual variance for the averaged endpoint becomes
\(\sigma^2 / n_{\text{rep}}\) while the random slope
contribution \(\sigma_b^2 t^2\) is essentially unchanged.
Similarly, the baseline visit at \(j = 0\) can be
replicated to reduce \(\sigma^2\) in the change-score
covariance.

\end{orig}

\begin{orig}
This is distinct from the common close design (Section
5): common close observations are at \emph{new} time
points after treatment stops, providing information
about the off-treatment slope. Replicated endpoints are
at the \emph{same} time point as the last treatment
visit, providing only measurement-error reduction.

\end{orig}

\subsection{Variance formula with replication}\label{variance-formula-with-replication}

\begin{orig}
The residual covariance matrix \(R\) is modified so that
the diagonal element corresponding to a replicated
visit uses \(\sigma^2 / n_{\text{rep}}\) instead of
\(\sigma^2\), and the off-diagonal elements reflecting
the shared baseline error use
\(\sigma^2 / n_{\text{rep,pre}}\). The Woodbury
computation proceeds identically.

\end{orig}

\begin{table}

\caption{\label{tab:replication-table}Fractional variance reduction from replicating the final endpoint ($J_0=0$, $J_1=2$, $J_2=0$).}
\centering
\begin{tabular}[t]{rrrrr}
\toprule
$n_{\text{rep}}$ & $\sigma_b/\sigma=0.5$ & $\sigma_b/\sigma=1$ & $\sigma_b/\sigma=2$ & $\sigma_b/\sigma=5$\\
\midrule
1 & 0.000 & 0.000 & 0.000 & 0.000\\
2 & 0.182 & 0.091 & 0.030 & 0.005\\
3 & 0.250 & 0.125 & 0.042 & 0.007\\
5 & 0.308 & 0.154 & 0.051 & 0.009\\
\bottomrule
\end{tabular}
\end{table}

\begin{orig}
The table confirms that endpoint replication is most
beneficial when measurement error dominates
(\(\sigma_b / \sigma\) small). When between-subject slope
variability is large (\(\sigma_b / \sigma > 2\)), the
gain from replication is modest because the variance
is dominated by \(\sigma_b^2 t^2\), which replication
cannot reduce.

\end{orig}

\subsection{Replication vs additional run-in observations}\label{replication-vs-additional-run-in-observations}

\begin{orig}
A natural question is whether a fixed budget of
additional measurements should be spent on run-in
observations (at distinct time points, separated by
\(t\)) or on endpoint replicates (closely spaced,
measurement-error reduction only).

\end{orig}

\begin{figure}
\centering
\pandocbounded{\includegraphics[keepaspectratio,alt={(\#fig:rep\_vs\_runin)Variance as a function of additional measurements allocated either to run-in (new time points) or endpoint replication (same time point). Baseline: J\_0=0, J\_1=2, \textbackslash sigma\^{}2=1, \textbackslash sigma\_b\^{}2=1.}]{report_files/report_files/figure-latex/rep_vs_runin-1.pdf}}
\caption{(\#fig:rep\_vs\_runin)Variance as a function of additional measurements allocated either to run-in (new time points) or endpoint replication (same time point). Baseline: \(J_0=0\), \(J_1=2\), \(\sigma^2=1\), \(\sigma_b^2=1\).}
\end{figure}

\begin{orig}
At \(\sigma_b / \sigma = 1\) (the setting of the figure),
the comparison depends on how many extra measurements
are available. With a single extra measurement, the
endpoint replicate is the better use (variance 2.73
versus 2.82 for one run-in observation): the lone
run-in observation pays the degree-of-freedom penalty
of Section 3. From the second extra measurement onward,
run-in observations at new time points win (2.16 versus
2.63 at two extras, 1.64 versus 2.57 at three), because
they provide information about the individual slope
\(b_i\), not just measurement-error reduction. When
\(\sigma_b / \sigma < 1\), endpoint replication is
stronger still: at \(\sigma_b / \sigma = 0.5\) it beat
run-in at every budget we examined (one to four extra
measurements), since measurement error is then the
dominant source of variance.

\end{orig}

\subsection{Combined replication at both ends}\label{combined-replication-at-both-ends}

\begin{orig}
Replicating both the baseline and endpoint visits
provides symmetric measurement-error reduction.

\end{orig}

\begin{table}

\caption{\label{tab:both-ends}Per-pair variance with baseline and endpoint replication ($J_0=0$, $J_1=2$, $\sigma^2=1$, $\sigma_b^2=1$).}
\centering
\begin{tabular}[t]{rrrr}
\toprule
$n_{\text{pre}}$ & $n_{\text{post}}=1$ & $n_{\text{post}}=2$ & $n_{\text{post}}=3$\\
\midrule
1 & 3.0000 & 2.7273 & 2.6250\\
2 & 2.7273 & 2.5000 & 2.4138\\
3 & 2.6250 & 2.4138 & 2.3333\\
\bottomrule
\end{tabular}
\end{table}

\section{Incorporating dropout}\label{incorporating-dropout}

\subsection{Pattern-mixture approach}\label{pattern-mixture-approach}

\begin{orig}
Following\textsuperscript{\citeproc{ref-frost2008}{13}} (Section 2.3), dropout
is handled via a pattern-mixture model\textsuperscript{\citeproc{ref-dawson1993}{15},\citeproc{ref-lu2008}{16}}. Participants
are stratified by their observed visit pattern, and
the overall variance is a weighted average:

\end{orig}

\begin{equation}
  N = \frac{1}{\sum_{s} p_s / N_s}
\end{equation}

\begin{orig}
where \(p_s\) is the proportion with pattern \(s\) and
\(N_s\) is the sample size that would be needed if all
participants had pattern \(s\). Under monotone dropout,
the pattern for a participant who leaves after
post-baseline visit \(j\) is the truncated design that
retains the run-in and baseline visits plus the first
\(j\) post-baseline observations; a participant with no
post-baseline visit contributes no information about
the treatment contrast, so that stratum enters with
\(N_s = \infty\). Dropout during the pre-randomization
run-in is ignored here, since run-in dropouts are
typically replaced before randomization. This
approximation treats each stratum independently.\textsuperscript{\citeproc{ref-hu2021}{17}} derive exact variance formulas under monotone
dropout for the three-component random coefficient
regression model (\(\sigma^2, \sigma_a^2, \sigma_b^2\)),
accounting for the effect of dropout on the
estimability of the random slope variance. Their
approach is more rigorous than the pattern-mixture
approximation used here, which is conservative when
dropout is informative about the random slope.

\end{orig}

\begin{table}

\caption{\label{tab:dropout}Sample size adjusted for dropout via pattern-mixture ($J_0=2$, $J_1=2$, $J_2=1$, MIRIAD parameters).}
\centering
\begin{tabular}[t]{rrr}
\toprule
Annual dropout & $N$ (no dropout) & $N$ (with dropout)\\
\midrule
0.0 & 104 & 104\\
0.1 & 104 & 128\\
0.2 & 104 & 160\\
0.3 & 104 & 204\\
\bottomrule
\end{tabular}
\end{table}

\section{Numerical examples}\label{numerical-examples}

\subsection{AD neuroimaging trial}\label{ad-neuroimaging-trial}

\begin{orig}
Using MIRIAD parameters (\(\sigma^2 = 0.3363\),
\(\sigma_b^2 = 0.9745\), \(t = 1\) year), compare the
optimal allocation for \(J = 4, 5, 6\) total
observations with and without dropout.

\end{orig}

\begin{table}

\caption{\label{tab:ad-optimal}Top 10 allocations for $J=5$ observations (MIRIAD parameters, $\Delta=0.25$, 90\% power).}
\centering
\begin{tabular}[t]{lrr}
\toprule
$(J_0, J_1, J_2)$ & $\mathop{\mathrm{Var}}_1(\hat\gamma)$ & $N$ per group\\
\midrule
(0,3,2) & 0.291783 & 50\\
(0,2,3) & 0.303230 & 51\\
(1,2,2) & 0.453050 & 77\\
(2,3,0) & 0.551570 & 93\\
(0,4,1) & 0.560024 & 95\\
\addlinespace
(1,3,1) & 0.587888 & 99\\
(3,2,0) & 0.593958 & 100\\
(2,2,1) & 0.616148 & 104\\
(0,1,4) & 0.701145 & 118\\
(1,1,3) & 0.803845 & 136\\
\bottomrule
\end{tabular}
\end{table}

\section{Discussion}\label{discussion}

\begin{orig}
The central allocation result is negative and, we believe,
useful: under a fixed total visit budget, the optimal design
places no observations in the run-in phase in any
configuration examined, including the Alzheimer's
neuroimaging regime (Table \ref{tab:optimal-grid}). The
budget splits between the treatment and common close phases.
The mechanism is the degree-of-freedom cost identified in
Section 3: any positive run-in count adds the
pre-randomization slope parameter \(\delta\), while run-in
observations inform the treatment contrast only indirectly.
This does not contradict the efficiency gains reported for
run-in designs in Paper 1 and the prior literature; those
gains arise when run-in visits are additional measurements
appended to a fixed treatment phase. When visits are
fungible, as under a per-participant visit budget, the same
visit buys more information in the treatment or common close
phase. Trial designers should therefore ask which question
their operational constraint poses: extra visits (run-in can
help) or reallocated visits (run-in never helps in the
settings examined here).

\end{orig}

\begin{orig}
The common close allocation is governed jointly by the
signal-to-noise ratio and the budget. A free common close
observation always helped in the configurations we examined
(Section 5), so the operative question is the budgeted
trade-off: moving one visit from the treatment phase to the
common close reduces the variance once \(\sigma_b / \sigma\)
exceeds a crossover that falls with the budget, from about
0.81 at \(J = 3\) to 0.30 at \(J = 8\)
(Table \ref{tab:common-close-crossover}). Above the
crossover the optimizer assigns the close phase a growing
share of the budget, up to an even split at \(J = 8\) in the
settings examined. In practice, high slope-to-noise outcomes
such as neuroimaging endpoints sit well above the crossover
at realistic budgets, so a substantial common close phase is
warranted there, while low slope-to-noise outcomes justify a
common close only at larger budgets. Estimating
\(\sigma_b / \sigma\) from pilot data is the prerequisite
design step.

\end{orig}

\begin{orig}
Two refinements operate within a fixed phase allocation.
First, when observation times within a phase are free, the
two-point raw-outcome special case of Section 6 indicates
that spacing observations toward the extremes of the window
minimizes the slope-estimation variance; the corresponding
optimization in the change-score model of this paper remains
open. Second, replicating the baseline or endpoint visit
reduces the measurement-error contribution to the change
score without adding new time points; this is most valuable
at low slope-to-noise ratios (up to a 31 percent variance
reduction at \(\sigma_b / \sigma = 0.5\) with five replicates)
and offers little when slope variability dominates (at most
5 percent at \(\sigma_b / \sigma = 2\)). The comparison of
Section 7 clarifies the choice between spending an
additional measurement on a new run-in time point versus an
endpoint replicate: at \(\sigma_b / \sigma = 1\) the replicate
wins when only one extra measurement is available, because a
lone run-in observation pays the degree-of-freedom penalty,
while run-in observations win from the second extra
measurement onward by contributing slope information; below
\(\sigma_b / \sigma = 1\) replication won at every budget
examined.

\end{orig}

\begin{orig}
Several assumptions bound the scope of these recommendations.
The allocation optima are computed under equally spaced
observations, linear individual trajectories, and a
compound-symmetric residual structure; the AR(1) analysis of
Paper 2 implies that under decaying correlation the optimum
shifts toward fewer and more closely spaced run-in
observations, since distant ones lose informativeness. The
allocation results themselves are numerical, obtained by
exhaustive enumeration of the feasible phase splits; an
analytical characterization of the optimum, for example by
signing the marginal value of a close-phase visit, remains
open. Dropout is accommodated through the pattern-mixture
approximation of Section 8, applied to monotone truncation
strata, which the exact monotone-dropout treatment of\textsuperscript{\citeproc{ref-hu2021}{17}} would refine. Extending the allocation optimization
jointly over phase counts, within-phase timing, and the
correlation structure remains the natural next step.

\end{orig}

\section{Conflict of interest}\label{conflict-of-interest}

\begin{orig}
The author declares no conflicts of interest.

\end{orig}

\section{Funding}\label{funding}

\begin{orig}
This work received no specific funding.

\end{orig}

\section{Data availability statement}\label{data-availability-statement}

\begin{orig}
The \texttt{runinpower} R package implementing the
optimal-allocation routines (\texttt{allocation.R} and
\texttt{replicated.R}, with the design matrices they rely on built
in \texttt{design.R}), validation tests, and the shared
bibliography across the four-paper compendium are openly
available in the project repository. Specific paths are
listed in the Reproducibility section below.

\end{orig}

\section{Reproducibility}\label{reproducibility}

\begin{orig}
This manuscript was rendered with R 4.4.x. Principal packages
used are the in-package \texttt{runinpower} source, \texttt{tinytest}
(testing harness), and \texttt{rmarkdown} plus \texttt{bookdown} for the
manuscript build.

\end{orig}

\begin{orig}
\textbf{Random-number generator.} All results in this manuscript
are deterministic evaluations of the variance formulas given
fixed input parameters; no Monte Carlo simulation is used
and no random seed is required.

\end{orig}

\begin{orig}
\textbf{Data and code availability.} The R package source is at
\texttt{R/}. The companion manuscripts (Papers 1, 2, 4) are at
\texttt{analysis/report/\{01,02,04\}-*/}. The manuscript source and
shared bibliography (\texttt{references.bib}, a symlink to
\texttt{../../references.bib}) are at
\texttt{analysis/report/03-allocation/}.

\end{orig}

\section{References}\label{references}

\protect\phantomsection\label{refs}
\begin{CSLReferences}{0}{1}
\bibitem[\citeproctext]{ref-winkens2005}
\CSLLeftMargin{1. }%
\CSLRightInline{Winkens B, Schouten HJA, Breukelen GJP van, Berger MPF. Optimal time-points in clinical trials with linearly divergent treatment effects. \emph{Statistics in Medicine}. 2005;24(24):3743-3756. doi:\href{https://doi.org/10.1002/sim.2370}{10.1002/sim.2370}}

\bibitem[\citeproctext]{ref-yi2002}
\CSLLeftMargin{2. }%
\CSLRightInline{Yi Q, Panzarella T. Estimating sample size for tests on trends across repeated measurements with missing data based on the interaction term in a mixed model. \emph{Controlled Clinical Trials}. 2002;23(5):481-496. doi:\href{https://doi.org/10.1016/S0197-2456(02)00228-3}{10.1016/S0197-2456(02)00228-3}}

\bibitem[\citeproctext]{ref-frison1992}
\CSLLeftMargin{3. }%
\CSLRightInline{Frison L, Pocock SJ. Repeated measures in clinical trials: Analysis using mean summary statistics and its implications for design. \emph{Statistics in Medicine}. 1992;11(13):1685-1704. doi:\href{https://doi.org/10.1002/sim.4780111304}{10.1002/sim.4780111304}}

\bibitem[\citeproctext]{ref-overall1994}
\CSLLeftMargin{4. }%
\CSLRightInline{Overall JE, Doyle SR. Estimating sample sizes for repeated measurement designs. \emph{Controlled Clinical Trials}. 1994;15(2):100-123. doi:\href{https://doi.org/10.1016/0197-2456(94)90015-9}{10.1016/0197-2456(94)90015-9}}

\bibitem[\citeproctext]{ref-galbraith2002}
\CSLLeftMargin{5. }%
\CSLRightInline{Galbraith S, Marschner IC. Guidelines for the design of clinical trials with longitudinal outcomes. \emph{Controlled Clinical Trials}. 2002;23(3):257-273. doi:\href{https://doi.org/10.1016/S0197-2456(02)00205-2}{10.1016/S0197-2456(02)00205-2}}

\bibitem[\citeproctext]{ref-atkinson1992}
\CSLLeftMargin{6. }%
\CSLRightInline{Atkinson AC, Donev AN. \emph{Optimum Experimental Designs}. Clarendon Press; 1992.}

\bibitem[\citeproctext]{ref-fedorov1997}
\CSLLeftMargin{7. }%
\CSLRightInline{Fedorov VV, Hackl P. \emph{Model-Oriented Design of Experiments}. Vol 125. Springer; 1997. doi:\href{https://doi.org/10.1007/978-1-4612-0703-0}{10.1007/978-1-4612-0703-0}}

\bibitem[\citeproctext]{ref-mentre1997}
\CSLLeftMargin{8. }%
\CSLRightInline{Mentré F, Mallet A, Baccar D. Optimal design in random-effects regression models. \emph{Biometrika}. 1997;84(2):429-442. doi:\href{https://doi.org/10.1093/biomet/84.2.429}{10.1093/biomet/84.2.429}}

\bibitem[\citeproctext]{ref-tekle2008}
\CSLLeftMargin{9. }%
\CSLRightInline{Tekle FB, Tan FES, Berger MPF. D-optimal cohort designs for linear mixed-effects models. \emph{Statistics in Medicine}. 2008;27(14):2586-2600. doi:\href{https://doi.org/10.1002/sim.3080}{10.1002/sim.3080}}

\bibitem[\citeproctext]{ref-nash2021}
\CSLLeftMargin{10. }%
\CSLRightInline{Nash S, Morgan KE, Frost C, Mulick A. Power and sample-size calculations for trials that compare slopes over time: Introducing the slopepower command. \emph{The Stata Journal}. 2021;21:572-601. doi:\href{https://doi.org/10.1177/1536867X211045512}{10.1177/1536867X211045512}}

\bibitem[\citeproctext]{ref-iddi2022}
\CSLLeftMargin{11. }%
\CSLRightInline{Iddi S, Donohue MC. Power and sample size for longitudinal models in {R} -- the longpower package and {Shiny} app. \emph{The R Journal}. 2022;14(1):264-282. \url{https://journal.r-project.org/articles/RJ-2022-022/}}

\bibitem[\citeproctext]{ref-winkens2007}
\CSLLeftMargin{12. }%
\CSLRightInline{Winkens B, Breukelen GJP van, Schouten HJA, Berger MPF. Randomized clinical trial designs for auto-correlated longitudinal data. \emph{Statistics in Medicine}. 2007;26(28):5023-5037. doi:\href{https://doi.org/10.1002/sim.2907}{10.1002/sim.2907}}

\bibitem[\citeproctext]{ref-frost2008}
\CSLLeftMargin{13. }%
\CSLRightInline{Frost C, Kenward MG, Fox NC. Optimizing the design of clinical trials where the outcome is a rate. {Can} estimating a baseline rate in a run-in period increase efficiency? \emph{Statistics in Medicine}. 2008;27(18):3717-3731. doi:\href{https://doi.org/10.1002/sim.3280}{10.1002/sim.3280}}

\bibitem[\citeproctext]{ref-goulet2019}
\CSLLeftMargin{14. }%
\CSLRightInline{Goulet MA, Cousineau D. The power of replicated measures to increase statistical power. \emph{Advances in Methods and Practices in Psychological Science}. 2019;2:199-213.}

\bibitem[\citeproctext]{ref-dawson1993}
\CSLLeftMargin{15. }%
\CSLRightInline{Dawson JD, Lagakos SW. Size and power of two-sample tests of repeated measures data. \emph{Biometrics}. 1993;49(4):1022-1032.}

\bibitem[\citeproctext]{ref-lu2008}
\CSLLeftMargin{16. }%
\CSLRightInline{Lu K, Luo X, Chen PY. Sample size estimation for repeated measures analysis in randomized clinical trials with missing data. \emph{International Journal of Biostatistics}. 2008;4(1):Article 9. doi:\href{https://doi.org/10.2202/1557-4679.1098}{10.2202/1557-4679.1098}}

\bibitem[\citeproctext]{ref-hu2021}
\CSLLeftMargin{17. }%
\CSLRightInline{Hu N, Mackey H, Thomas R. Power and sample size for random coefficient regression models in randomized experiments with monotone missing data. \emph{Biometrical Journal}. 2021;63(4):806-824. doi:\href{https://doi.org/10.1002/bimj.202000184}{10.1002/bimj.202000184}}

\end{CSLReferences}

\end{document}
